\section{Related work}
\label{sec:related}

\paragraph{Evolution and learning.}
That learning can guide evolution without any inheritance of acquired traits was shown in simulation by \citet{hinton1987}; \citet{fontanari2017} argue that those simulations demonstrated the selection of learning ability rather than the Baldwin effect proper, and estimate the probability and time with which a learned behaviour becomes innate. \citet{ackley1992} combined evolution and learning in an artificial environment of adaptive and non-adaptive organisms and report that the combination was more successful than either alone; \citet{le2019baldwin} report a similar result for self-taught neural networks in a multi-agent foraging and survival simulation. \citet{roupassov2026} provide an open-source reimplementation of a 1992 artificial-life testbed for evolutionary reinforcement learning and compare neural with programmatic policies using survival analysis over 4,000 trials. Our individuals resemble programmatic policies (a fixed rule skeleton with a few evolved parameters), and our setting adds a third source of competence that none of these studies includes: a frozen pretrained model whose judgements are fixed and whose contribution can be switched off, shuffled or replaced.

\paragraph{Language models as priors for agents.}
\citet{yan2025llmprior} treat a language model as a prior over actions and integrate it into reinforcement learning through Bayesian inference, reporting a large reduction in required samples in offline settings. \citet{arumugam2025} have language models implement posterior sampling to improve exploration. \citet{murugesan2020} use commonsense knowledge from ConceptNet to let agents in text-based games filter implausible actions. These studies ask whether a prior speeds up learning; they do not separate the contribution of an evolved genome from that of learning, and they do not test the conditions under which a prior is neutral or unnecessary. A different line uses language models inside evolutionary search: \citet{lehman2022} use a code model as a mutation operator, and \citet{sygkounas2026} co-evolve language-model-generated policies and environments. In our work the model is not part of the search; it supplies fixed judgements that evolved genomes learn how much to trust. The frozen model we use, \laya{}, is described by its model card as a base to be specialised rather than a zero-shot decision engine \citep{laya}; we report, under our setting and with a correction for its preference for the first option, how well its zero-shot judgements of food descriptions separate nutritious from toxic ones (Section~\ref{sec:res-quality}). Its encoder is ModernBERT \citep{warner2024}.

\paragraph{Delayed consequences.}
In animals, an aversion to a food can be learned even when illness follows long after eating, and the cue that becomes associated depends on the kind of consequence \citep{garcia1966}; \citet{kalat1973} proposed that the long delay is bridged by ``learned safety'': during the delay the animal learns that the familiar taste is safe. Our credit-window gene, which attributes a delayed shock to foods eaten shortly before, is an engineering analogue of the problem these studies describe; we make no claim that our individuals model the biological mechanism.

\paragraph{Evaluation practice.}
Pre-registration separates predictions from post hoc analyses \citep{nosek2018}. In reinforcement learning the number of random seeds determines the risk of statistical errors \citep{colas2018}; we treat the training seed as the unit of analysis, state effect-size thresholds in advance, and use evolution strategies \citep{salimans2017} for training. As a sentiment baseline we use the VADER lexicon \citep{hutto2014}.
